% TOP_PROBLEM: {"id": 4700016, "title": "Reversible equivariant planar differential systems", "category_id": 11, "category": {"id": 11, "name": "dynamical_systems", "display_name": "Dynamical Systems", "description": "Problems about long-term behavior of deterministic systems, Hamiltonian dynamics, and periodic orbits.", "slug": "dynamical-systems", "order_index": 11, "created_at": "2026-07-31T15:26:25.671Z"}, "status": "partially_solved"}
% FIRST_POSTED: null
% ENTRY_KIND: statement_only
% Imported from: batch15.tex; UnsolvedMath ID 4700016
\documentclass[11pt]{article}
\usepackage[margin=25mm]{geometry}
\usepackage{fontspec}
\setmainfont{Latin Modern Roman}
\usepackage{amsmath,amssymb,mathtools}
\usepackage{unicode-math}
\setmathfont{Latin Modern Math}
\usepackage{xurl,hyperref}
\hypersetup{colorlinks=true,urlcolor=blue}
\setcounter{secnumdepth}{0}
\setlength{\parindent}{0pt}
\setlength{\parskip}{5pt}
\setlength{\emergencystretch}{3em}
\newcommand{\sref}[2]{\hyperlink{source:#1:#2}{[S#2]}}
\newcommand{\cataloguescope}[1]{\paragraph{Recorded target.}#1}
\begin{document}
% UnsolvedMath ID 4700016; source catalogue code AMR-046-0016
\setcounter{section}{4700015}
\section{Reversible equivariant planar differential systems}\label{4700016}
{\small\sffamily UnsolvedMath ID 4700016 · Dynamical Systems}
\subsection{Definitions and mathematical statement}

\paragraph{Shared notation.}$\mathbb N=\{1,2,\ldots\}$, and $\mathbb Z,\mathbb Q,\mathbb R,\mathbb C$ denote the integers, rationals, reals and complex numbers. Any other conventions are specified locally.


Identify the plane with $\mathbb C$ by $z=x+iy$, and fix integers $n,k\ge1$. Consider the real autonomous differential equation \[\dot z=iz+(z\bar z)^n z^{k+1}.\] Its reversible center at the origin is surrounded by a period annulus, the maximal connected annular region filled with periodic orbits that shrink toward the origin. Parameterize these orbits by their intersection $z=r>0$ with the positive real ray, ordered from the center outward. Write $T(r)$ for the minimal positive period in the original time variable. Monotone decrease means $T(r_2)\le T(r_1)$ for any $r_1<r_2$ in this annulus.
\paragraph{Statement recorded in the supplied source.}
For every pair $n,k\ge1$, is the period function $T$ monotonically decreasing throughout the period annulus of the origin? \sref{4700016}{2}
\subsection{Short English statement}
Does the oscillation period decrease from the center outward for every member of this reversible complex polynomial family?
\subsection{Sources}
\begin{description}
\item[\hypertarget{source:4700016:1}{S1}] ULAM.AI and the UnsolvedMath Contributors, \emph{UnsolvedMath: A Curated Collection of Open Mathematics Problems}, record 4700016 (AMR-046-0016). CC BY 4.0. \url{https://huggingface.co/datasets/ulamai/UnsolvedMath}.
\item[\hypertarget{source:4700016:2}{S2}] A. Gasull, Some open problems in low dimensional dynamical systems (2020), Section 3 introductory period definitions; Section 3.5, equation (15), Proposition 3.1 and Problem 16, pp.13,15--16. \url{https://arxiv.org/pdf/2012.02524}.
\end{description}
\label{end:4700016}
\end{document}
